%============================================================
\chapter{Bai Hoc: Thong Minh Va Su Sang Tao (Intelligence and Creativity)}
\begin{center}
  \textit{\color{truyenblue}Intelligence is not knowing many things but knowing how to use what you know.}\\
  \textit{(Khon ngoan khong phai la biet nhieu dieu ma la biet su dung dieu minh biet.)}\\[4pt]
  \small\color{truyengold}--- Thien nhien day ta ---
\end{center}
\ngancach

\section{Qua Vo Va Con Quang}
\begin{truyen}{Qua Vo Va Con Quang}{Con Quang}
\chuhoa{M}{}M ot con quang va tinh co nhin thay cai qua oc nat gan ben duong nhung khong the be vo.\\
\textit{(A clever crow found a walnut near the road but could not crack it open.)}

No bay len cao, thu tha qua oc xuong mat duong cung ma van khong be duoc.\\
\textit{(It flew high and dropped the walnut onto the hard road but still could not crack it.)}

Sau do no sang kien the: no cho xe o to chay qua nghi qua oc se bi nghien nat.\\
\textit{(Then it had a clever idea: it waited for a car to drive over, thinking the nut would be crushed.)}

Nhung no phai cho sang den do de xe dung lai roi no moi gat lay nhan khong bi xe chet.\\
\textit{(But it had to wait for the red light so the car stopped and then it safely collected the kernel.)}

Con quang da ket hop vat ly, giao thong va thoi gian de giai quyet mot bai toan don gian.\\
\textit{(The crow combined physics, traffic, and timing to solve a simple problem.)}

\end{truyen}
\begin{baihoc}
  Su thong minh that su la tim giai phap dac biet cho van de binh thuong, khong phai lam phuc tap them.\\
  \textit{(True intelligence is finding a unique solution for an ordinary problem, not making it more complicated.)}
\end{baihoc}

\section{Bot Bong Ton Den Va Hoa}
\begin{truyen}{Bot Bong Ton Den Va Hoa}{Ong Bong}
\chuhoa{K}{}K hi mua dong den, dong vat muon giu am bang nhieu cach khac nhau dac biet.\\
\textit{(When winter arrives, animals stay warm in many special different ways.)}

Dan ong vo dat thanh qua cau lon o trung tam to de thu nhiet the phat ra.\\
\textit{(The bee colony forms a large ball at the center of the hive to capture body heat.)}

Nhung con o ngoai cung run that nhanh bam chat vao nhau de tao nhiet ma khong phai di chuyen.\\
\textit{(Those on the outer shell vibrate fast and cling tightly together generating heat without moving.)}

Sau do chung luan doi vi tri: con lanh vao trong, con am quay ra ngoai thay phien.\\
\textit{(Then they rotate positions: cold ones move inside, warm ones move outside in turns.)}

He thong luan doi do giu toan dan on am suot mua dong ma khong can nhien lieu bao nhieu.\\
\textit{(That rotation system keeps the whole colony warm through winter using very little fuel.)}

\end{truyen}
\begin{baihoc}
  Thong minh tap the thuong viet hon bat ky tri tue ca nhan nao: hay tan dung moi nao luc cua nhom.\\
  \textit{(Collective intelligence surpasses any individual intellect: harness all the brainpower of your team.)}
\end{baihoc}

\section{Bo Ho Tach Mu Con Moi}
\begin{truyen}{Bo Ho Tach Mu Con Moi}{Bo Ho}
\chuhoa{B}{}B o ho can pha cac vet thuong tren da bo de xua duoi cac con rui ruoi muon de trung.\\
\textit{(Oxpecker birds pick wounds on buffalo skin to drive away flies that want to lay eggs.)}

Nhung dieu bat ngo la bo ho cung an luon sau bo dang an da bo tren vet thuong.\\
\textit{(But surprisingly, oxpeckers also eat the lice currently eating buffalo hide on the wounds.)}

Bo ho len tren lung trau bo khong chi an sau bo ma con dong vai tro bao ve canh gac.\\
\textit{(Oxpeckers on buffalo backs not only eat parasites but also serve as sentinel guards.)}

Khi co su tu den gan, bo ho keu to canh bao cho trau bo chay thi ngay.\\
\textit{(When a lion approaches, the oxpecker cries out loudly warning the buffalo to run immediately.)}

Moi quan he ky sinh thoat hoat thanh moi quan he cong sinh khi hai ben biet nhin xa hon.\\
\textit{(A parasitic relationship transforms into a mutualistic one when both sides know to look farther ahead.)}

\end{truyen}
\begin{baihoc}
  Thong minh nhat la bien nguoi tung la ke lay loi ich tu ban thanh nguoi dong hanh giup ich cho ban.\\
  \textit{(The smartest move is turning someone who once extracted benefit from you into a partner who benefits you.)}
\end{baihoc}

\section{Ca Tam Giac Luoi San}
\begin{truyen}{Ca Tam Giac Luoi San}{Ca Tam Giac}
\chuhoa{L}{}L oan ca tam giac lon song o vung nuoc am luon san moi bang chien thuat nguoc sang.\\
\textit{(Schools of large triggerfish living in warm waters always hunt prey by moving against the light.)}

Chung thuong tan cong khi mat troi o sau lung chung de mat troi chieu vao mat con moi.\\
\textit{(They often attack when the sun is behind them so sunlight shines directly into the prey's eyes.)}

Con moi bi loa mat boi anh sang choi chang khong kip nhan biet nguy hiem sap den.\\
\textit{(The prey is blinded by the glare and unable to recognize the approaching danger in time.)}

Ca tam giac su dung vat ly quang hoc phu hop voi thoi gian trong ngay mot cach chinh xac.\\
\textit{(Triggerfish use optical physics aligned with the exact time of day with precision.)}

Cac nha quan su hoc chien thuat tan cong nguoc sang tu loai ca nay de ung dung vao chien tranh.\\
\textit{(Military strategists learned the counter-sun attack tactic from this fish to apply to warfare.)}

\end{truyen}
\begin{baihoc}
  Biet su dung moi truong xung quanh nhu vu khi la dinh cao cua tri tue chien luoc.\\
  \textit{(Knowing how to use the surrounding environment as a weapon is the peak of strategic intelligence.)}
\end{baihoc}

\section{Voi Dung Cong Cu}
\begin{truyen}{Voi Dung Cong Cu}{Voi}
\chuhoa{N}{}N guoi ta tung nghi chi con nguoi moi biet su dung cong cu de thuc hien cong viec.\\
\textit{(People once thought only humans knew how to use tools to perform tasks.)}

Nhung voi biet cat canh cay va cam vao voi de quat ruoi va muoi mot cach hieu qua.\\
\textit{(But elephants know to pick up branches and hold them with their trunks to swat flies and mosquitoes effectively.)}

Voi cung biet lay da to che vao mieng lo nuoc de nuoc khong bay hoi trong ngay nang nong.\\
\textit{(Elephants also know to use large stones to cover water holes so water doesn't evaporate on hot days.)}

Mot so con con biet dung thanh cay dai xoa len tuong de ve hinh theo y thich rieng.\\
\textit{(Some even use long sticks to draw on walls creating pictures according to their own preferences.)}

Tri tue khong goi han o bat ky loai nao; no hien huu o noi nao co nhu cau va kham pha.\\
\textit{(Intelligence is not limited to any species; it exists wherever there is need and curiosity.)}

\end{truyen}
\begin{baihoc}
  Hay mo rong dinh nghia cua ban ve tri tue; nhung giai phap tot nhat co the den tu noi ban it ky vong nhat.\\
  \textit{(Broaden your definition of intelligence; the best solutions may come from where you least expect them.)}
\end{baihoc}

\section{Khi Dau Tien Dung Gguong}
\begin{truyen}{Khi Dau Tien Dung Gguong}{Con Khi}
\chuhoa{L}{}L an dau tien nguoi ta dat guong truoc chuong khi, tat ca bong doi mat o la thay.\\
\textit{(The first time researchers placed a mirror before a monkey cage, all stared in confusion at the stranger.)}

Nhungo mot so loai khi cao cap nhu tinh tinh sau vai phut hieu ra minh dang nhin chinh minh.\\
\textit{(But some higher primates like chimpanzees realized after a few minutes they were looking at themselves.)}

Tinh tinh bat dau kiem tra mat, rang, long nguoi ban gom bam duoi mat guong.\\
\textit{(The chimpanzee began examining its face, teeth, and fur that was stuck under its eyes using the mirror.)}

No dung guong nhu cong cu de lam dep va kiem tra khu vuc co the no khong tu nhin thay duoc.\\
\textit{(It used the mirror as a tool for grooming and checking body areas it cannot see directly.)}

Tu nhan thuc la buoc dau tien cua tri tue cao cap ma rat it sinh vat tren Trai Dat co duoc.\\
\textit{(Self-awareness is the first step of higher intelligence that very few creatures on Earth possess.)}

\end{truyen}
\begin{baihoc}
  Bat dau hieu ban than la cuoc hanh trinh thong minh nhat ma bat ky ai co the bat tay vao.\\
  \textit{(Beginning to understand yourself is the most intelligent journey anyone can embark upon.)}
\end{baihoc}

\section{Nuoc Dang Len Den Ca Tho}
\begin{truyen}{Nuoc Dang Len Den Ca Tho}{Chim Co}
\chuhoa{M}{}M ot con chim co khong the chen mo vao binh co co mieng nho de uong nuoc.\\
\textit{(A crow could not fit its beak into a bottle with a narrow mouth to drink water.)}

No nhin xung quanh va bat dau nhat da nho ne trong vao binh nuoc chiem dat.\\
\textit{(It looked around and began picking up small pebbles to drop into the water bottle occupying space.)}

Muc nuoc tu tu dang len theo tung hon da no them vao mot cach kien nhan.\\
\textit{(The water level slowly rose with each pebble it patiently added one by one.)}

Sau bay muoi hon da, nuoc len du cao de no thoe mo uong duoc no ne no long.\\
\textit{(After seventy pebbles, the water rose high enough for it to dip its beak in and drink happily.)}

Bay gio bai toan nay duoc day trong truong tieu hoc nhu vi du ve tu duy sang tao.\\
\textit{(Today this problem is taught in elementary school as an example of creative thinking.)}

\end{truyen}
\begin{baihoc}
  Khi khong the di theo duong thang, hay tim cach gian tiep them den de vat can hoa thanh loi the.\\
  \textit{(When you cannot go straight, find an indirect way to add to the obstacle until it becomes an advantage.)}
\end{baihoc}

\section{Ca Voi Dung Am Thanh San Moi}
\begin{truyen}{Ca Voi Dung Am Thanh San Moi}{Ca Voi Xanh}
\chuhoa{A}{}A m thanh cua ca voi xanh co the truyen di xa den 1600 km trong dat nuoc truoc khi tan bien.\\
\textit{(The sound of a blue whale can travel up to 1,600 kilometers through ocean water before disappearing.)}

Chung su dung am thanh de tao ban do tinh thon sac nu xac ca hang nghin km vuong xung quanh.\\
\textit{(They use sound to create detailed topographic maps of the seabed thousands of square kilometers around them.)}

Khi loai ca moi boi tu xuong sau len mat nuoc, ca voi da do vi tri chinh xac truoc do 20 phut.\\
\textit{(When prey fish swim up from the deep to the surface, whales have pinpointed their exact location 20 minutes before.)}

Ca voi dung vat ly am thanh trong nuoc nhu he thong GPS song am tu tu nhieu chieu.\\
\textit{(Whales use underwater acoustic physics like a multidimensional sonar GPS system.)}

Chung da thong minh hon nhieu loai cong cu ma con nguoi moi phat minh ra gan day.\\
\textit{(They were smarter than many tools humans only recently invented.)}

\end{truyen}
\begin{baihoc}
  Tri tue that su khong can nhan mat; no co the lang nghe, cam nhan va ve ban do the gioi bang phuong tien khac.\\
  \textit{(True intelligence needs no eyes; it can listen, sense, and map the world by other means.)}
\end{baihoc}

\section{Nhen Day Cuu Minh}
\begin{truyen}{Nhen Day Cuu Minh}{Nhen}
\chuhoa{M}{}M ot con nhen roi xuong ao nuoc va khong the boi vao bo duoc vi dong nuoc khong ung ho.\\
\textit{(A spider fell into a pond and could not swim to shore because the water current did not help.)}

No tuc thi dung chan sau de xoay than tao ra luc day nguoc chieu dong nuoc.\\
\textit{(It immediately used its back legs to spin its body creating a force against the current direction.)}

Dong nuoc day no vao bo thay vi ra giua ho theo nhu phuong phap benh vien nguoc.\\
\textit{(The current pushed it toward shore instead of toward the center of the lake, like a reverse paddle method.)}

Toan bo hanh dong chi mat vai giay nhung la ket qua cua viec hieu vat ly chat long.\\
\textit{(The whole action took only seconds but was the result of understanding fluid physics.)}

Khoa hoc vat ly khong chi la sach giao khoa; no la ky nang song cua moi sinh vat tren dot.\\
\textit{(Physics is not just a textbook; it is the survival skill of every living creature on Earth.)}

\end{truyen}
\begin{baihoc}
  Hoc vat ly, hoa hoc, sinh hoc khong phai de thi ma de song; thien nhien da biet dieu do tu lau.\\
  \textit{(Study physics, chemistry, and biology not to pass tests but to live; nature has always known this.)}
\end{baihoc}

\section{Biet Doi Du Bao Thoi Tiet}
\begin{truyen}{Biet Doi Du Bao Thoi Tiet}{Biet Doi}
\chuhoa{K}{}K hi ap suat khi quyen thay doi truoc con bao, con nguoi kho nhan ra song biet doi tuc thi cam nhan duoc.\\
\textit{(When atmospheric pressure changes before a storm, humans barely notice but bats immediately sense it.)}

Biet doi bay thap hem, co lai long va tim noi truy nen hin an nang xuat doc tuc thoi.\\
\textit{(Bats fly lower, tuck their wings, and find sheltered spots as their instinctive forecast activates.)}

Cac nha khoa hoc theo doi hanh vi bat thuong cua biet doi de du bao bao to chinh xac hon.\\
\textit{(Scientists track abnormal bat behavior to forecast storms more accurately.)}

Co the nhay cam cua loai vat thuong vuot xa bat ky may moc nao con nguoi che tao ra.\\
\textit{(Animals' sensitive bodies often far surpass any instruments humans have ever created.)}

Thong minh khong chi la phan tich so lieu ma con la biet lang nghe tin hieu thien nhien.\\
\textit{(Intelligence is not just analyzing data but also knowing how to listen to nature's signals.)}

\end{truyen}
\begin{baihoc}
  Hanh vi cua dong vat quanh ban co the la ban tin thoi tiet chinh xac hon bat ky ung dung dien thoai nao.\\
  \textit{(Animal behavior around you can be a more accurate weather report than any phone app.)}
\end{baihoc}
